\chapter{Astrometry: Distances, Motions, and the Space Velocity Vector}

The coordinate systems of Chapter~2 fix a star's \emph{direction} on the sky but say nothing about its distance or its motion through space. Astrometry is the branch of astrophysics dedicated to precise measurements of the positions, distances, and kinematics of celestial bodies. By combining angular coordinates $(\alpha, \delta)$, trigonometric parallax ($p$), proper motions $(\mu_\alpha^*, \mu_\delta)$, and spectroscopic radial velocities ($v_r$), astronomers reconstruct the full six-dimensional phase-space distribution of stars in the Galaxy: position $\vecr = (x, y, z)$ and space velocity $\vec v = (v_x, v_y, v_z)$. In this chapter, we develop the foundations of astrometric measurements, distance determination, and 3D velocity reconstruction.

\section{Trigonometric Parallax and the Parsec}

As the Earth orbits the Sun, a nearby star appears to trace a tiny elliptical path on the sky relative to distant background objects, for exactly the same reason that a nearby lamppost appears to shift against a distant skyline as you walk sideways past it. This apparent angular shift is the \emph{annual trigonometric parallax}.

\begin{definition}[Trigonometric Parallax and Parsec]
	Consider the right triangle formed by the Sun, the Earth, and a star at distance $d$, at the instant Earth's orbital position gives the maximum angular offset (when the Earth--Sun line is perpendicular to the line of sight, i.e.\ the Sun itself would appear $90^\circ$ from the star as seen from Earth). This triangle has a right angle at the Sun, an adjacent side of length $1\,\mathrm{AU}$ (the Earth--Sun distance), a hypotenuse $d$ (the Sun--star distance), and an angle $p$ at the star between the hypotenuse and the line to the Sun. Elementary trigonometry on this right triangle gives
	\begin{equation}
		\sin p = \frac{1\,\mathrm{AU}}{d} \quad\Longrightarrow\quad p \approx \frac{1\,\mathrm{AU}}{d}\quad (\text{for } p \ll 1\,\text{rad}),
		\label{eq:parallax_def}
	\end{equation}
	where $p$ (or $\varpi$), the \emph{parallax angle}, is by construction the semi-major axis of the star's apparent parallactic ellipse on the sky over the year, and $1\,\mathrm{AU} = 1.495978707\times10^{11}\,\mathrm{m}$ is the Earth--Sun distance.
	A \emph{parsec} ($\mathrm{pc}$) is defined as the distance at which an astronomical unit subtends an angle of exactly one arcsecond ($1''$):
	\begin{keyresult}
		\begin{equation}
			d\,[\mathrm{pc}] = \frac{1}{p\,[\mathrm{arcsec}]}.
			\label{eq:parallax_distance_pc}
		\end{equation}
	\end{keyresult}
	In metric units:
	\begin{equation}
		1\,\mathrm{pc} = \frac{1\,\mathrm{AU}}{\tan(1'')} = \frac{1.4959787 \times 10^{11}\,\mathrm{m}}{\pi / (180 \times 3600)} \approx 3.08567758 \times 10^{16}\,\mathrm{m} \approx 3.26156\,\mathrm{ly}.
	\end{equation}
\end{definition}

\begin{example}[Distance to Sirius]
	Sirius has a measured parallax $p = 379.21\,\mathrm{mas} = 0.37921''$ (from the \emph{Hipparcos}/\emph{Gaia} catalogs). Its distance is:
	\[
	d = \frac{1}{0.37921} \approx 2.637\,\mathrm{pc} \approx 8.60\,\mathrm{ly}.
	\]
\end{example}

\paragraph{Parallax Inversion and the Lutz-Kelker Bias.}
When converting a measured parallax $p_{\rm obs}$ into distance, non-linear error propagation means $\langle 1/p \rangle \neq 1/\langle p \rangle$. For fractional parallax errors $\sigma_p/p \gtrsim 0.15$, symmetrical Gaussian errors on $p$ produce strongly asymmetric, biased distance estimates (the \emph{Lutz-Kelker bias}), requiring Bayesian distance priors (e.g., exponentially decreasing space density models in \emph{Gaia} data releases).

\section{Three-Dimensional Cartesian Position Vector}

From the distance $d$ and equatorial coordinates $(\alpha, \delta)$, the heliocentric Cartesian position vector is:
\begin{keyresult}
	\begin{equation}
		\vecr = \begin{pmatrix} x \\ y \\ z \end{pmatrix} = d\,\hat{\mathbf r}_{\rm eq} = \begin{pmatrix} d\cos\delta\cos\alpha \\ d\cos\delta\sin\alpha \\ d\sin\delta \end{pmatrix}.
		\label{eq:cartesian_pos_3d}
	\end{equation}
\end{keyresult}

\section{Proper Motion and Transverse Velocity}

Stars possess intrinsic space velocities relative to the Sun. The component of stellar motion perpendicular to the line of sight projects onto the celestial sphere as an angular displacement over time, known as \emph{proper motion} $\boldsymbol{\mu}$.

\begin{definition}[Proper Motion Components]
	Proper motion is resolved into two orthogonal components on the sky:
	\begin{itemize}
		\item $\mu_\delta = \dfrac{d\delta}{dt}$: Rate of change of declination ($\mathrm{mas/yr}$ or $\mathrm{arcsec/yr}$).
		\item $\mu_\alpha^* \equiv \mu_\alpha\cos\delta = \dfrac{d\alpha}{dt}\cos\delta$: True angular rate of change along the parallel of declination ($\mathrm{mas/yr}$).
	\end{itemize}
	The total proper motion magnitude is:
	\begin{equation}
		\mu = \sqrt{(\mu_\alpha\cos\delta)^2 + \mu_\delta^2} = \sqrt{(\mu_\alpha^*)^2 + \mu_\delta^2}.
		\label{eq:total_proper_motion}
	\end{equation}
\end{definition}

\paragraph{Derivation of the $4.74047$ Kinematic Conversion Constant.}
The physical transverse velocity $v_t$ (in $\mathrm{km/s}$) across a distance $d$ (in $\mathrm{pc}$) with proper motion $\mu$ (in $\mathrm{arcsec/yr}$) is:
\[
v_t = d\,\mu = \left(d\,[\mathrm{pc}] \times 3.085678 \times 10^{13}\,\mathrm{km/pc}\right) \times \left(\mu\,[\mathrm{arcsec/yr}] \times \frac{1\,\mathrm{rad}}{206264.806''}\right) \times \left(\frac{1\,\mathrm{yr}}{3.15576 \times 10^7\,\mathrm{s}}\right).
\]
Evaluating the numerical coefficient:
\[
k \equiv \frac{3.085678 \times 10^{13}}{206264.806 \times 3.15576 \times 10^7} \approx 4.7404705\,\mathrm{km\cdot s^{-1}\cdot yr\cdot arcsec^{-1}\cdot pc^{-1}}.
\]
\begin{keyresult}
	\begin{equation}
		v_t\,[\mathrm{km/s}] = 4.74047 \times \mu\,[\mathrm{arcsec/yr}] \times d\,[\mathrm{pc}] = 4.74047 \times \frac{\mu\,[\mathrm{arcsec/yr}]}{p\,[\mathrm{arcsec}]}.
		\label{eq:vt_master}
	\end{equation}
\end{keyresult}

\begin{example}[Transverse Velocity of Sirius]
	Sirius has proper motion $\mu_\alpha^* = -546.01\,\mathrm{mas/yr} = -0.54601''/\mathrm{yr}$ and $\mu_\delta = -1223.07\,\mathrm{mas/yr} = -1.22307''/\mathrm{yr}$, giving total proper motion:
	\[
	\mu = \sqrt{(-0.54601)^2 + (-1.22307)^2} \approx 1.3396''/\mathrm{yr}.
	\]
	With $d = 2.637\,\mathrm{pc}$, its transverse velocity is:
	\[
	v_t = 4.74047 \times 1.3396 \times 2.637 \approx 16.75\,\mathrm{km/s}.
	\]
\end{example}

\section{Radial Velocity and the Doppler Effect}

The component of stellar velocity along the line of sight, $v_r$, is measured spectroscopically via the Doppler shift of absorption or emission lines.

\begin{definition}[Doppler Shift and Radial Velocity]
	Consider a source emitting successive wave crests spaced by one period $T_0 = \lambda_0/c$ in its own rest frame, receding directly away from the observer at speed $v_r \ll c$. Between the emission of one crest and the next, the source recedes an extra distance $v_r T_0$, so the second crest must travel that much farther to reach the observer, arriving a time $v_r T_0/c$ later than it would from a stationary source. The observed period is therefore stretched to $T_{\rm obs} = T_0 + v_rT_0/c = T_0(1+v_r/c)$, and since $\lambda = cT$, the relative wavelength shift $z$ is:
	\begin{keyresult}
		\begin{equation}
			z \equiv \frac{\Delta\lambda}{\lambda_0} = \frac{\lambda_{\rm obs} - \lambda_0}{\lambda_0} = \frac{T_{\rm obs}-T_0}{T_0} = \frac{v_r}{c},
			\label{eq:doppler_classical}
		\end{equation}
	\end{keyresult}
	valid for non-relativistic velocities ($v_r \ll c$), where $v_r > 0$ denotes recession (redshift) and $v_r < 0$ denotes approach (blueshift).
\end{definition}

For relativistic velocities (cosmological redshift, relativistic jets), the exact special-relativistic Doppler relation must be used:
\begin{equation}
	1 + z = \sqrt{\frac{1 + v_r/c}{1 - v_r/c}} \quad\Longrightarrow\quad \frac{v_r}{c} = \frac{(1+z)^2 - 1}{(1+z)^2 + 1}.
	\label{eq:doppler_relativistic}
\end{equation}

\section{The Full 3D Space Velocity Vector}

Combining proper motion components and radial velocity, the space velocity vector in the spherical equatorial basis $\{\hat{\mathbf r}, \hat{\boldsymbol{\alpha}}, \hat{\boldsymbol{\delta}}\}$ is:
\begin{equation}
	\vec v = v_r\,\hat{\mathbf r} + v_\alpha\,\hat{\boldsymbol{\alpha}} + v_\delta\,\hat{\boldsymbol{\delta}},
\end{equation}
where $v_\alpha = 4.74047\,\mu_\alpha^*\,d$ and $v_\delta = 4.74047\,\mu_\delta\,d$.

The directions $\hat{\boldsymbol\alpha}$ and $\hat{\boldsymbol\delta}$ needed here are simply the local compass directions ``increasing $\alpha$'' and ``increasing $\delta$'' tangent to the celestial sphere at the star's position --- precisely the (normalized) partial derivatives of the position unit vector $\hat{\mathbf r}$ (Eq.~\eqref{eq:unit_vector_eq}) with respect to each spherical coordinate, a standard construction for a local tangent basis on a sphere:
\begin{equation}
	\hat{\boldsymbol\delta} = \pdv{\hat{\mathbf r}}{\delta} = \begin{pmatrix} -\sin\delta\cos\alpha \\ -\sin\delta\sin\alpha \\ \cos\delta \end{pmatrix}, \qquad
	\hat{\boldsymbol\alpha} = \frac{1}{\cos\delta}\pdv{\hat{\mathbf r}}{\alpha} = \begin{pmatrix} -\sin\alpha \\ \cos\alpha \\ 0 \end{pmatrix},
\end{equation}
where the $1/\cos\delta$ in $\hat{\boldsymbol\alpha}$ restores unit length ($\partial\hat{\mathbf r}/\partial\alpha$ alone has magnitude $\cos\delta$, since displacements in $\alpha$ shrink physically toward the poles). One can check directly that $\hat{\mathbf r}$, $\hat{\boldsymbol\alpha}$, $\hat{\boldsymbol\delta}$ are mutually orthogonal unit vectors at every point, so together they form the orthonormal equatorial coordinate basis:
\begin{equation}
	\hat{\mathbf r} = \begin{pmatrix} \cos\delta\cos\alpha \\ \cos\delta\sin\alpha \\ \sin\delta \end{pmatrix}, \qquad
	\hat{\boldsymbol{\alpha}} = \begin{pmatrix} -\sin\alpha \\ \cos\alpha \\ 0 \end{pmatrix}, \qquad
	\hat{\boldsymbol{\delta}} = \begin{pmatrix} -\sin\delta\cos\alpha \\ -\sin\delta\sin\alpha \\ \cos\delta \end{pmatrix}.
	\label{eq:equatorial_basis_vectors}
\end{equation}

Transforming to the equatorial Cartesian velocity frame $\vec v = (v_x, v_y, v_z)^T$:
\begin{keyresult}
	\begin{equation}
		\begin{pmatrix} v_x \\ v_y \\ v_z \end{pmatrix} = \begin{pmatrix} \cos\delta\cos\alpha & -\sin\alpha & -\sin\delta\cos\alpha \\ \cos\delta\sin\alpha & \cos\alpha & -\sin\delta\sin\alpha \\ \sin\delta & 0 & \cos\delta \end{pmatrix} \begin{pmatrix} v_r \\ v_\alpha \\ v_\delta \end{pmatrix}.
		\label{eq:cartesian_velocity_transform}
	\end{equation}
\end{keyresult}
The total speed through space relative to the Sun is:
\begin{equation}
	v_{\rm total} = \sqrt{v_r^2 + v_t^2} = \sqrt{v_x^2 + v_y^2 + v_z^2}.
	\label{eq:v_total_speed}
\end{equation}

\section{Galactic Space Velocity Components $(U, V, W)$}

In Galactic dynamics, stellar velocities are expressed in a right-handed Cartesian frame centered on the Sun:
\begin{itemize}
	\item $U$: Velocity directed toward the Galactic Center ($l = 0^\circ, b = 0^\circ$).
	\item $V$: Velocity directed along the direction of Galactic rotation ($l = 90^\circ, b = 0^\circ$).
	\item $W$: Velocity directed toward the North Galactic Pole ($b = +90^\circ$).
\end{itemize}

Because $(U,V,W)$ and $(v_x,v_y,v_z)$ are both Cartesian components of the \emph{same} physical velocity vector, expressed in two frames sharing a common origin (the Sun), they are related by a pure rotation --- the same rigid-rotation argument used in Chapter~2 to transform coordinates between celestial frames, applied here to a velocity rather than a position. The Galactic velocity vector $(U, V, W)^T$ is obtained from the equatorial Cartesian velocity $(v_x, v_y, v_z)^T$ via the transformation matrix $\mathbf T_{\rm Gal}$, whose numerical entries are fixed once and for all by the IAU-adopted direction of the Galactic pole (Chapter~2, Eq.~\eqref{eq:gal_b}):
\begin{keyresult}
	\begin{equation}
		\begin{pmatrix} U \\ V \\ W \end{pmatrix} = \mathbf T_{\rm Gal} \begin{pmatrix} v_x \\ v_y \\ v_z \end{pmatrix} = \begin{pmatrix} -0.054876 & -0.873437 & -0.483835 \\ +0.494109 & -0.444830 & +0.746982 \\ -0.867666 & -0.198076 & +0.455984 \end{pmatrix} \begin{pmatrix} v_x \\ v_y \\ v_z \end{pmatrix}.
		\label{eq:uvw_transform}
	\end{equation}
\end{keyresult}

\paragraph{Correction to the Local Standard of Rest (LSR).}
The Local Standard of Rest is a reference frame moving on a circular orbit around the Galactic Center at the Sun's Galactocentric radius ($R_0 \approx 8.2\,\mathrm{kpc}$) with circular speed $v_0 \approx 220\text{--}240\,\mathrm{km/s}$. The Sun's peculiar motion with respect to the LSR (Schönrich et al. 2010) is:
\begin{equation}
	(U_\odot, V_\odot, W_\odot) = (11.10^{+0.69}_{-0.75},\, 12.24^{+0.47}_{-0.47},\, 7.25^{+0.37}_{-0.36})\,\mathrm{km/s}.
\end{equation}
The stellar velocity relative to the LSR is:
\begin{equation}
	(U_{\rm LSR}, V_{\rm LSR}, W_{\rm LSR}) = (U + U_\odot, V + V_\odot, W + W_\odot).
\end{equation}

\section{The Cosmic Distance Ladder}

Trigonometric parallax provides the foundational, geometric rung of the cosmic distance ladder without astrophysical assumptions. Beyond the reach of direct parallax ($d \gtrsim \text{few kpc}$ from ground, $\sim 10\text{--}20\,\mathrm{kpc}$ from \emph{Gaia}), successive secondary and tertiary distance indicators are calibrated:

\begin{enumerate}
	\item \textbf{Spectroscopic Parallax:} Inferring absolute magnitude $M_V$ from spectral type and luminosity class on the H-R diagram, determining distance from $m - M = 5\log_{10}d - 5$.
	\item \textbf{Main-Sequence Fitting:} Matching the color-magnitude diagram of star clusters to the zero-age main sequence (ZAMS).
	\item \textbf{Standard Candles (Cepheids and RR Lyrae):} Leavitt's Period-Luminosity relation for Classical Cepheid variables ($M_V \approx -2.81\log_{10}P - 1.43$) bridges distances out to $\sim 30\,\mathrm{Mpc}$ (Hubble Space Telescope key project).
	\item \textbf{Type Ia Supernovae:} White dwarf thermonuclear detonations standardized via the Phillips light-curve width-luminosity relation ($M_B \approx -19.3$), extending distances to cosmological scales ($z \sim 1\text{--}2$).
	\item \textbf{The Hubble-Lemaître Law:} On cosmological scales, galaxy recession velocities follow the expansion of spacetime:
	\begin{equation}
		v_r = H_0\,d, \qquad H_0 \approx 70\,\mathrm{km\cdot s^{-1}\cdot Mpc^{-1}}.
		\label{eq:hubble_law}
	\end{equation}
\end{enumerate}

Every rung of this distance ladder is calibrated photometrically -- through an apparent brightness, a color, or a light-curve shape -- which is precisely the subject developed next in Chapter~4.
